%%%%%%%%%%%%%%%%%%%%%%%%%%%%%%%%%%%%%%%%%
% Seismica
% LuaLaTeX Template, version 2025
%
% Important note:
% This template must be compiled with LuaLaTeX, the below lines will ensure this
%!TEX TS-program = lualatex
%!TEX encoding = UTF-8 Unicode
%%%%%%%%%%%%%%%%%%%%%%%%%%%%%%%%%%%%%%%%%


% Options available: report, breakmath, proof, onecolumn, invited, opinion, languages
% breakmath breaks the equations that are too long
% proof prints a proof watermark
% onecolumn only if issue with breakmath
% invited prints "invited" in the banner and headers
% languages loads additional languages to babel. English is set as the default main language.

\documentclass[breakmath,report,languages={french,japanese,arabic}]{seismica}
%\documentclass[breakmath,opinion,proof]{seismica}

%% If report, specify the type to appear in header (comment or uncomment). if not specified, will only show 'report'
%\reporttype{Fast Report}
%\reporttype{Null Results Report}
%\reporttype{Software Report}
%\reporttype{Data Report}

% Languages
% Specific fonts need to be imported for languages that do not use latin alphabet.
% For each language, you can have a look at the babel doc (https://texdoc.org/serve/babel/0, section 1.7) and click on the reference mark (※) to open a specific site for each language with examples.
\babelfont[japanese]{rm}{Harano Aji Mincho}
\babelfont[japanese]{sf}{Harano Aji Mincho}
\babelfont[arabic]{rm}{FreeSerif}
\babelfont[arabic]{sf}{FreeSerif}

% Metadata by SCE team only
\dois{10.26443/seismica.v2i1.N}
\receiveddate{September 23, 2022}
\accepteddate{January 6, 2023}
\publisheddate{January 11, 2023}
\theyear{2023}
\thevolume{2}
\thenumber{1}  %% issue number
\prodedname{Efirstname Editorname}  %% Production editor
\handedname{Efirstname Editorname}  %% handling editor
\copyedname{Cfirstname Copyedname}  %% copyed editor
%\translatorname{translator name, if there is one}
%\reviewername{Rfirstname Reviewname\\ Rfirstname Reviewname}  % if at least one non-anonymous

%%%%%%%%%%%%%%%%%%%%%%%%%%%%%%%%%%%%%%%%%
% Metadata by authors (Replace everything below)

\title{A template for Seismica}
\shorttitle{A template for Seismica} % used for header

% use ~ for non-breaking spaces in authors, so names don't break across lines
\author[1]{A. Author1
	\thanks{Corresponding author: bla@som.ac.edu}
	\orcid{0000-0002-1825-0097}}
\author[1]{B.~author2 
	\orcid{0000-0002-1825-097}}
\author[2]{C.~Author3
	\orcid{0000-0002-1825-007}}
\affil[1]{affil Author 1 and 2 }
\affil[2]{affil author 3}

\credit{Funding acquisition}{Alice, Bob}
\credit{Writing}{Charlie, Doris}
\credit{Writing - review \& editing}{Emilio, Francis}


\begin{document}

% English asbtract is first
\makeseistitle{%
	\begin{summary}{Abstract}
	Replace text
	\vspace{1.5cm} % because abstract is too short
	\end{summary}
	\begin{summary}{Non-technical summary}
	The text goes here. Again, no longer than 200 words.
	\vspace{1.5cm} % because abstract is too short
	\end{summary}
}

 %%%% In case there are too many abstracts and need a page break, use this:
\addsummaries{
    \begin{selectlanguage}{japanese} 
        \begin{summary}{要約}
        日本語
            従来の短期地震活動モデル（例えば時間的疫病型余震列（ETAS）モデル）は、地震カタログデータのみに依存し、地殻変形などの外部強制メカニズムを組み込んでいない。本研究では、GNSS測定から得られたひずみ速度情報を統合することでETASモデルを拡張する
        \end{summary}
    \end{selectlanguage}
    \begin{selectlanguage}{french} 
        \begin{summary}{Résumé} 
        Texte en francais
        \end{summary}
    \end{selectlanguage}
    \begin{selectlanguage}{arabic} 
        \begin{summary}{ملخص} 
        تعتمد نماذج النشاط الزلزالي القصير المدى الكلاسيكية، مثل نموذج تسلسل الهزات الارتدادية من النوع الوبائي (ETAS)، بشكل حصري على بيانات كتالوج الزلازل ولا تدمج آليات التأثير الخارجية مثل تشوه القشرة الأرضية. في هذه الدراسة، قمنا بتوسيع نموذج ETAS من خلال دمج معلومات معدل الإجهاد المستمدة من قياسات GNSS، مما يسمح بتغير المعدل الأساسي بمرور الوقت من خلال تداخل قائم على البيانات مع نواة استجابة مقدرة تجريبياً.
        \end{summary}
    \end{selectlanguage}
}  %% don't forget this one!
	
	\section{Introduction}
	
	Cite with \citep{metropolis_monte_1949} or \citet{metropolis_monte_1949}. \citet{van_den_ende_creating_2021} is an example of a preprint citation that can be adapted for other unpublished work. \citet{rf_software} is an example of a software citation that can be adapted for other code, datasets, seismic networks, etc. \citet{mueller_documentation_2003} is an example of a report citation, which can be adapted for technical and other reports.
	
	To refer to a figure, use Fig.~\ref{fig:2} or Figs~\ref{fig:1}, \ref{fig:2} (Tab. and Tabs).
	
	\begin{figure}[ht!]
		\includegraphics[width=\columnwidth]{empty} 
		\caption{column-wide figure.}
		\label{fig:1}
	\end{figure}
	
	\section{Section 1}
	
	\subsection{Subsection }
	
	\begin{figure*}[ht!]
		\includegraphics[width=\textwidth]{empty} 
		\caption{Full-width figure.}
		\label{fig:2}
	\end{figure*}

Refer to equation \ref{eq1}.

\begin{equation} \label{eq1}
\mathbf{G} = \frac{1}{2}(2\cos z) + (1/2)(2\cos z+j\sin z-j\sin z) + (1/2)(\cos z+j\sin z+\cos z-j\sin z) -  (1/2)(e^{jz}+e^{-jz})
\end{equation}
	
\begin{table}[ht!]

\caption{In tables, prefer to use \code{m}-type columns, as they wrap the text. Specify the width of each column in brackets: \code{m{width}}. You can use \code{\textbackslash hline} or \code{\textbackslash thickhline} for horizontal rules. Also, you can add a \code{\textbackslash rule{0pt}{2ex}} after \code{hlines} to add an extra vertical space.} 
\label{tab:etas_parameters}

\begin{tabular}{ m{2.5cm} m{2.5cm} m{2.5cm}  } % m-type columns allow for text wrapping, specify the width in brackeds
\textbf{Parameter} & \textbf{Stationary Model} & \textbf{Deformation-Driven Model} \\
\hline
\( \mu_0 \) [events/day] & 0.0084 & 0.026 \\
\( K_0 \)                & 0.057  & 0.055 \\
\( c \) [days]           & 2.8 \(\times 10^{-4}\) & 3.0 \(\times 10^{-4}\) \\
\( \alpha \)             & 0.47   & 0.48 \\
\( p \)                  & 0.94   & 0.95 \\
AIC                      & 638    & 634 \\
\end{tabular}
\end{table}

	\subsection{Code}

Code examples should be concise and descriptive. They should introduce core functionality or specific syntax and should be included using the \code{lstlisting} environment. Note that lines longer than 45 characters will be broken when using the prepress option. Extended examples or use cases should be uploaded separately. Individual words of code can be written inline, for example:

To improve stability of the inversion, the \code{Model} object accepts the \code{strict} keyword, which disables piecewise linear approximation of the target function (Listing~\ref{code}).

\begin{lstlisting}[caption=Example use of \code{Model}, label=code, language=Python]
#2 4 6 8 0 2 4 6 8 0 2 4 6 8 0 2 4 6 8 0 2 4|
import mymodule as mm

model = mm.Model(strict=True)
mdls = model.perturb()

for mdl in mdls:
	var = mdl.get_variance()
\end{lstlisting}

\section{Page formatting}
If you need to add a landscape-format page, for a table or figure, there are header/footer styles in the template. Adding a landscape page will mess with tex's page-filling algorithm since it will insert the rotated page exactly where it appears in the text, so you will need to manually move text before/after to maximize page fill.

%% add these lines (uncommented) before the landscape content
%\onecolumn
%\thispagestyle{lscapedplain}
%\pagestyle{lscapedplain}
%\begin{landscape}

%% [add your landscape content here - if a figure*, use height=\textwidth (or .9\textwidth) instead of width to set size.

%% add these lines (uncommented) after the landscape content to reset to portrait mode
%\end{landscape}
%\restoregeometry
%\thispagestyle{plainReset}
%\pagestyle{plainReset}

% If you want a caption for a full-page landscape figure to be on the next page but still have the same number, you can use
%\addcounter{figure}{-1}
% before adding another figure environment with just the caption following the landscape page. Note that this will not play well with the XML conversion.


\begin{acknowledgements}
	Thank all relevant parties and acknowledge funding sources, if any.
\end{acknowledgements}
	
\section*{Data availability}
Authors should direct readers to an open access repository such as figshare or Github, where data are made available.

\section*{Competing interests}
The authors declare no competing interests.

\bibliography{biblio} 
	
\end{document}

